\paragraph{Proof of intrinsic non-degeneracy.}
The identity below uses the metric on \(\mathcal S(\theta)\) induced by the
ambient Euclidean space; accordingly, \(I_k\) denotes the identity in any
orthonormal basis of
\(T_{x^\star(\theta)}\mathcal S(\theta)\).  Fix \(\theta\), abbreviate
\[
R=R(\theta)=\sqrt{1+\theta^2},
\qquad
b_\theta=1+0.3\theta,
\]
and write \(x=(Rv,0)\) with \(v\in\mathbb S^k\).  On the minimizer
manifold, \(\lVert x\rVert=R\) and
\(u(\theta,x)=\langle v,a(\theta)\rangle\), so the restricted upper objective
is
\[
\left.f(\theta,\cdot)\right|_{\mathcal S(\theta)}(Rv,0)
=2\sqrt{\epsilon^2+R^2}
+q_\theta\!\left(\langle v,a(\theta)\rangle\right).
\]
The first term is constant on the manifold.

Let
\(\eta=(\eta_y,0)\in T_{x^\star(\theta)}\mathcal S(\theta)\), set
\(r=\lVert\eta\rVert\), and, for \(r>0\), let \(w=\eta_y/r\).  Since
\(x^\star(\theta)=(-Ra(\theta),0)\), we have
\(\langle w,a(\theta)\rangle=0\).  The geodesic with
\(x_\eta(0)=x^\star(\theta)\) and \(\dot x_\eta(0)=\eta\) is
\[
x_\eta(t)
=\left(
R\left[
-a(\theta)\cos\!\left(\frac{rt}{R}\right)
+w\sin\!\left(\frac{rt}{R}\right)
\right],0
\right).
\]
Consequently,
\[
u(\theta,x_\eta(t))
=-\cos\!\left(\frac{rt}{R}\right),
\qquad
u'(0)=0,
\qquad
u''(0)=\frac{r^2}{R^2}.
\]
For
\[
q_\theta(u)
=b_\theta(u+1)
+\frac12\sin^2\!\left(2\pi(u+1)\right),
\]
the oscillatory term has zero first derivative at \(u=-1\), and hence
\(q_\theta'(-1)=b_\theta\).  The second-derivative chain rule now gives
\[
\begin{aligned}
\left\langle
\eta,\operatorname{Hess}_{\mathcal S(\theta)}
f(x^\star(\theta))\,\eta
\right\rangle
&=
\left.\frac{d^2}{dt^2}f(\theta,x_\eta(t))\right|_{t=0}\\
&=
q_\theta''(-1)\bigl(u'(0)\bigr)^2
+q_\theta'(-1)u''(0)\\
&=
\frac{1+0.3\theta}{1+\theta^2}\lVert\eta\rVert^2.
\end{aligned}
\]
Because this holds for every tangent vector \(\eta\),
\[
\operatorname{Hess}_{\mathcal S(\theta)}f(x^\star(\theta))
=\frac{1+0.3\theta}{1+\theta^2}I_k.
\]
Finally, for \(\theta\in[-3,3]\),
\[
\begin{aligned}
1+0.3\theta-0.01(1+\theta^2)
&=0.99+0.3\theta-0.01\theta^2\\
&\geq
\min\{0.99-0.9-0.09,\;0.99+0.9-0.09\}=0,
\end{aligned}
\]
where the minimum is attained at an endpoint because the quadratic is
concave.  Therefore
\[
\operatorname{Hess}_{\mathcal S(\theta)}f(x^\star(\theta))
\succeq0.01I_k,
\]
with equality in the lower bound at \(\theta=-3\).
\hfill\(\square\)
